\chapter{Coverage and Reference Appendix}

This appendix summarizes the current \api{crystal-qt6} API shape at a high level. Treat the README and specs as the executable source of truth when release boundaries move.

\begin{longtable}{@{}p{0.26\textwidth}p{0.66\textwidth}@{}}
\toprule
Area & Representative APIs \\
\midrule
Application & \api{Application}, \api{Qt6.application}, metadata, style sheets, clipboard, local event loops \\
Styles & \api{Style}, \api{CommonStyle}, \api{StyleFactory}, \api{ProxyStyle}, \api{StylePainter}, \api{StyledItemDelegate}, item editor factories, and the shared \api{StyleOption*} and \api{StyleHintReturn*} families \\
Shells & \api{MainWindow}, \api{Action}, \api{ActionGroup}, \api{MenuBar}, \api{Menu}, \api{ToolBar}, \api{StatusBar}, \api{DockWidget}, \api{SystemTrayIcon}, action icons, dock toggle/title-bar helpers, tray activation and notification flows, and status-bar size-grip control \\
Dialogs & \api{Dialog}, \api{MessageBox}, \api{FileDialog}, \api{ColorDialog}, \api{FontDialog}, \api{InputDialog}, \api{ProgressDialog}, \api{SplashScreen}, \api{DialogButtonBox}, aligned splash messages, and button-box layout helpers \\
Widgets & Labels, buttons, line edits, completers, check boxes, radio buttons, combo boxes, sliders, spin boxes, tabs, tool boxes, splitters, stacked widgets, scroll areas, group boxes, progress bars, frames, focus frames, LCD numbers, size grips, date/time controls, calendar widgets, text editors, text browsers, item widgets, and most of the ordinary widget/control surface the project currently targets \\
Events & \api{EventWidget}, paint, resize, mouse, wheel, key, focus, enter/leave, drag/drop callbacks, synthetic event helpers \\
Rendering & \api{QPainter}, \api{QPainterPath}, \api{QTransform}, \api{QPen}, \api{QBrush}, gradients, \api{QPalette}, \api{QRegion}, polygons, \api{QFont}, font metrics, SVG, and PDF \\
Graphics View & \api{GraphicsScene}, \api{GraphicsView}, \api{GraphicsItem}, graphics shape, text, pixmap, and group items, \api{GraphicsWidget}, proxy widgets, graphics layouts, transforms, effects, and scene events \\
Images & \api{QImage}, \api{QPixmap}, \api{QBitmap}, \api{QImageReader}, \api{QImageWriter}, raw buffers, pixel operations, conversion, transforms, scaling, and clipboard images \\
Undo & \api{UndoStack}, \api{UndoCommand}, clean state, macros, undo/redo actions, stack signals \\
Model/View & Item widgets, abstract list/tree models, standard item models, standard-item icons, proxy sorting/filtering, selection models, delegates, list/tree/table views, richer header interaction \\
\bottomrule
\end{longtable}

\noindent
For the narrative overview of the scene/view stack, item model, transforms, and scene-local panels, see Chapter 8. The appendix material below is meant to be the compact lookup companion to that chapter.

\section{Styles and Delegates Reference}

Use Chapter 12 for the conceptual walkthrough. This section is the shorter lookup map for the style-facing wrappers that now sit underneath application polish, offscreen control rendering, and item-view customization.

\begin{longtable}{@{}p{0.28\textwidth}p{0.64\textwidth}@{}}
\toprule
Type & Summary \\
\midrule
\api{Style} & Base wrapper for an active Qt style, including name lookup, standard palette access, and widget/application/palette polish hooks. \\
\api{CommonStyle} & Standalone common style instance when you need a concrete style object rather than just the application's current one. \\
\api{StyleFactory} & Static helper for listing available style keys and creating style instances such as Fusion. \\
\api{ProxyStyle} & Base-style wrapper and replacement surface for advanced style composition. It is useful infrastructure even though full proxy override hooks are not exposed yet. \\
\api{StylePainter} & Style-aware painter for drawing standard controls, primitives, complex controls, item text, and pixmaps onto widgets or offscreen paint devices. \\
\api{StyleOption} & Shared base packet for control geometry, state, palette, layout direction, and widget association. \\
\api{StyleOptionButton}, \api{StyleOptionProgressBar}, \api{StyleOptionSlider}, \api{StyleOptionTab}, \api{StyleOptionToolButton} & Control-specific option wrappers for low-level, style-aware paint and size state. \\
\api{StyleOptionViewItem} & Item-view row and cell paint state used by delegates and style-aware view rendering. \\
\api{StyledItemDelegate} & Style-aware model/view delegate for display formatting, paint hooks, size hints, editor creation, and editor commit flows. \\
\api{QItemEditorFactory}, \api{QItemEditorCreator}, \api{QStandardItemEditorCreator} & Shared editor registration surface for delegates that should create the same widgets for the same kinds of values. \\
\api{StyleHintReturn}, \api{StyleHintReturnMask}, \api{StyleHintReturnVariant} & Low-level style-hint result payloads used when a style query needs to return more than a simple integer or boolean. \\
\bottomrule
\end{longtable}

\section{Widget Reference}

The following table collects the major visible widget wrappers in one place. It is intentionally brief: use it as a map of the UI surface, then jump to the earlier chapters or the source when a widget becomes central to your application. The binding now covers the common widget and container families broadly enough that this appendix is mostly a navigation aid, not a list of speculative future wrappers.

\begin{longtable}{@{}p{0.28\textwidth}p{0.64\textwidth}@{}}
\toprule
Widget & Summary \\
\midrule
\api{Widget} & Generic \api{QWidget} wrapper for visibility, focus, sizing, palette, masking, repaint, and styling. It is the base for most visible controls. \\
\api{MainWindow} & Top-level application shell with central widget, menu bar, toolbars, status bar, and dock areas. Use it for desktop-style applications. \\
\api{Dialog} & Modal or modeless secondary window for settings, confirmation flows, and custom data-entry forms. \\
\api{MessageBox} & Standard confirmation, warning, error, and information dialog with button-result helpers. \\
\api{FileDialog} & Native-feeling file and directory picker for open, save, and import/export workflows. \\
\api{ColorDialog} & Color picker for accents, fills, layer colors, and alpha-aware visual settings. \\
\api{FontDialog} & Font picker for label styling, editor fonts, and typography choices. \\
\api{InputDialog} & Compact dialog for one-off text, integer, floating-point, or choice input. \\
\api{ProgressDialog} & Cancelable progress window for longer-running interactive tasks. \\
\api{SplashScreen} & Temporary startup window for branding and early status messages while the application loads. \\
\api{DialogButtonBox} & Standard dialog command row that keeps OK, Cancel, Help, and similar buttons aligned with Qt conventions. \\
\api{Label} & Read-only text or pixmap display for captions, status values, previews, and lightweight visual output. \\
\api{PushButton} & General command button for direct actions such as Apply, Save, Add, or Export. \\
\api{ToolButton} & Compact command button suited to toolbars, palettes, and menu-backed command launchers. \\
\api{CommandLinkButton} & Larger action button with descriptive secondary text, often used in launchers or task-oriented dialogs. \\
\api{CheckBox} & Independent boolean or tri-state option control. \\
\api{RadioButton} & Single-choice option inside a visible group of mutually related modes or settings. \\
\api{LineEdit} & Single-line text editor for names, paths, searches, and short freeform values. \\
\api{ComboBox} & Compact choice control that can also become editable for typed values with suggested options. \\
\api{FontComboBox} & Specialized combo box for selecting installed fonts. \\
\api{TextEdit} & Rich text editor with document and cursor access for formatted editing surfaces. \\
\api{PlainTextEdit} & Multiline plain-text editor for notes, logs, source snippets, and larger raw text blocks. \\
\api{TextBrowser} & Read-only rich-text viewer with HTML support and link handling. \\
\api{Slider} & Linear continuous-ish value control for zoom, volume, thresholds, and live-preview adjustments. \\
\api{Dial} & Rotary numeric control for compact dashboards, knobs, and instrument-like settings. \\
\api{ScrollBar} & Explicit scrolling control that can be used standalone or as part of a larger view shell. \\
\api{SpinBox} & Integer editor with bounded range and step controls. \\
\api{DoubleSpinBox} & Floating-point editor with range, precision, prefix/suffix, and step controls. \\
\api{DateTimeEdit} & Combined date/time editor for timestamp-style values. \\
\api{DateEdit} & Date-only editor for calendar values and schedule metadata. \\
\api{TimeEdit} & Time-only editor for clocks, durations, and schedule entry. \\
\api{CalendarWidget} & Graphical month calendar for date selection and schedule-oriented UIs. \\
\api{ProgressBar} & Read-only bounded value display for completion, loading, and background task progress. \\
\api{LcdNumber} & Segmented numeric display for dashboard counters and instrument-style readouts. \\
\api{Frame} & Decorative or structural frame used for separators, boxed panels, and low-level shell polish. \\
\api{FocusFrame} & Focus highlight widget that tracks another widget and lets Qt draw the platform-style focus indicator. \\
\api{GroupBox} & Titled panel container, optionally checkable, for visually grouping related settings. \\
\api{ToolBox} & Stacked-page inspector widget that shows one labeled section at a time. \\
\api{TabWidget} & Multi-page container with a built-in tab strip and a current-page API. \\
\api{TabBar} & Lightweight tab strip when page management lives elsewhere. \\
\api{StackedWidget} & Page stack without visible tabs; useful when another control chooses the active page. \\
\api{ScrollArea} & Single-child scrolling container for large forms, previews, and inspector panes. \\
\api{Splitter} & Resizable multi-pane container for editor shells and side-by-side work areas. \\
\api{DockWidget} & Repositionable side panel for layers, inspectors, logs, search results, and tool palettes. \\
\api{Wizard} and \api{WizardPage} & Page-based guided workflow surface for installers, setup flows, and multi-step task sequences with shared field registration. \\
\api{MdiArea} and \api{MdiSubWindow} & Multi-document workspace widgets for applications that keep several child editors or document panes alive at once. \\
\api{ListWidget} & Small list surface that owns its items directly. Use it for sidebars and simple collections. \\
\api{TreeWidget} & Hierarchical item widget for expandable owned data such as grouped layers or outline trees. \\
\api{TableWidget} & Grid widget that stores its own cell items for small editable tables. \\
\api{UndoView} & Read-only history list for \api{UndoStack} or \api{UndoGroup} state, useful in desktop editor shells. \\
\api{EventWidget} & Custom widget surface with direct paint, input, focus, and drag/drop callbacks. \\
\api{SizeGrip} & Explicit resize-corner widget for dialogs and utility windows. \\
\bottomrule
\end{longtable}

\section{Painting and Rendering Reference}

The painting surface is broader than the widget list because it spans value types, raster targets, vector geometry, theme-aware colors, and export helpers. This table is a quick orientation pass for the classes most likely to appear in custom widget painting, rendering examples, and export flows.

\begin{longtable}{@{}p{0.28\textwidth}p{0.64\textwidth}@{}}
\toprule
Type & Summary \\
\midrule
\api{Color} & RGBA color value used throughout painting, palettes, brushes, pens, and pixel operations. \\
\api{Point}, \api{PointF}, \api{LineF}, \api{Size}, \api{SizeF}, \api{Rect}, \api{RectF} & Common geometry value types for painter coordinates, image sizes, bounds, and layout-independent drawing math. \\
\api{QPainter} & Central drawing object for lines, shapes, text, images, paths, clipping, transforms, opacity, and composition onto a target. \\
\api{QPen} & Stroke description for outlines, including width, style, dash pattern, caps, joins, and color. \\
\api{QBrush} & Fill description for solid colors, brush styles, gradients, image textures, and pixmap textures. \\
\api{QPalette} & Theme-aware color source for application- or widget-level painting that should follow the active platform style. \\
\api{QPainterPath} & Vector path container for lines, curves, rectangles, ellipses, polygons, containment checks, and fill-rule-controlled geometry. \\
\api{QPainterPathStroker} & Utility that turns a stroked outline into a filled path, useful for hit testing and halos around thin geometry. \\
\api{QTransform} & Affine transform helper for translation, scaling, rotation, and mapping painter or path coordinates. \\
\api{QPolygon} & Integer-coordinate polygon for region construction, painter polygon drawing, and mask-style geometry. \\
\api{QPolygonF} & Floating-point polygon for painter polylines, polygons, and precise vector geometry work. \\
\api{QRegion} & Integer filled-area type for clip regions, widget masks, set operations, and region-scoped repaint scheduling. \\
\api{QLinearGradient} & Linear gradient fill source with stops, spread, and coordinate-mode control. \\
\api{QConicalGradient} & Angular gradient fill source for dials, badges, and radial color sweeps. \\
\api{QRadialGradient} & Radial gradient fill source with center, radius, focal point, stop inspection, and spread control. \\
\api{QImage} & General-purpose raster image with pixel read/write access, copy/scale/transform helpers, and a convenient offscreen paint target. \\
\api{QPixmap} & Display-oriented raster image that is good for widget grabs, cached UI artwork, and painter draw sources. \\
\api{QBitmap} & Monochrome pixmap variant used for masks, classic texture workflows, and region construction. \\
\api{QImageReader} & Configurable image decoder for probing metadata and loading raster formats from files or devices. \\
\api{QImageWriter} & Configurable image encoder for writing raster output with explicit format and quality control. \\
\api{QSvgRenderer} & SVG loader and renderer for drawing complete documents or named elements into an existing painter. \\
\api{QSvgWidget} & Widget-backed SVG display surface that is convenient for embedding vector art directly in the UI. \\
\api{QSvgGenerator} & Painter target that writes SVG output while reusing the same drawing code used for raster or widget painting. \\
\api{QPdfWriter} & Painter target for fixed-page PDF export, page sizing, and print-style document output. \\
\api{QFont} & Typeface, size, weight, and style description for painter text output and widget typography. \\
\api{QFontMetrics} & Integer text metrics for measuring width, ascent, descent, and layout spacing before drawing text. \\
\api{QFontMetricsF} & Floating-point text metrics that fit painter-coordinate work more naturally than integer metrics. \\
\bottomrule
\end{longtable}

\section{Graphics View Reference}

The Graphics View layer now has enough surface area to deserve its own appendix map. Use Chapter 8 for the mental model and workflows; use this table when you want a quick reminder of which wrapper to reach for.

\begin{longtable}{@{}p{0.28\textwidth}p{0.64\textwidth}@{}}
\toprule
Type & Summary \\
\midrule
\api{GraphicsScene} & Scene container for item ownership, scene bounds, background and foreground brushes, selection, focus, rendering, and scene-space item queries. \\
\api{GraphicsView} & Scrollable widget viewport over a scene, with scene attachment, item lookup, render hints, drag modes, transforms, mapping helpers, and view-side invalidation. \\
\api{GraphicsItem} & Base scene item with position, transforms, flags, selection, opacity, stacking, grouping, and scene-space geometry. \\
\api{GraphicsObject} & QObject-backed graphics item for signal-style property change callbacks, gesture hooks, event filters, and graphics effects. \\
\api{GraphicsWidget} & Layout-capable scene item with size policies, contents margins, focus policy, style, window-title helpers, and other panel-like widget behavior inside the scene. \\
\api{GraphicsProxyWidget} & Bridge for embedding ordinary \api{QWidget} controls inside the scene while keeping them addressable as scene items. \\
\api{GraphicsLayout} & Shared base for graphics-native layouts, including activation, invalidation, contents margins, size policies, and geometry queries. \\
\api{GraphicsLinearLayout} & Row or column layout for \api{GraphicsWidget} trees, nested layouts, stretch spacers, and per-item alignment or spacing. \\
\api{GraphicsGridLayout} & Grid layout for scene widgets with row and column spacing, alignment, stretch, and cell lookup. \\
\api{GraphicsAnchorLayout} and \api{GraphicsAnchor} & Constraint-style layout surface for edge and corner anchoring between graphics widgets. \\
\api{GraphicsRectItem}, \api{GraphicsEllipseItem}, \api{GraphicsLineItem} & Lightweight shape items with geometry access plus the familiar pen and brush styling surface. \\
\api{GraphicsPathItem} and \api{GraphicsPolygonItem} & Vector-geometry items for arbitrary paths, routes, outlines, and polygonal regions. \\
\api{GraphicsPixmapItem} & Pixmap-backed scene item for icons, markers, tokens, thumbnails, and raster overlays. \\
\api{GraphicsSimpleTextItem} and \api{GraphicsTextItem} & Scene text surfaces ranging from lightweight string rendering to document-backed rich text with cursor and link support. \\
\api{GraphicsItemGroup} & Grouping helper for moving or treating several scene items as a single unit while preserving item identities. \\
\api{GraphicsTransform}, \api{GraphicsRotation}, \api{GraphicsScale} & Reusable transform objects for ordered item transform lists beyond a single local rotation or scale value. \\
\api{GraphicsBlurEffect}, \api{GraphicsColorizeEffect}, \api{GraphicsDropShadowEffect}, \api{GraphicsOpacityEffect} & Item effects for emphasis, depth, tinting, softening, and overlay presentation. \\
\api{GraphicsSceneEvent} and subclasses & Scene-specific event wrappers for mouse, hover, wheel, drag/drop, context-menu, help, and move event inspection or synthetic test setup. \\
\bottomrule
\end{longtable}

\section{Enum Reference}

The following tables collect enum values that appear often in examples and callbacks. They are a reader-facing snapshot of the Crystal wrappers, not a substitute for checking the source when adding new bindings.

Flags enums can be combined with Crystal's bitwise OR operator. For example, an item can be made selectable and enabled with \api{Qt6::ItemFlag::Selectable | Qt6::ItemFlag::Enabled}.

\subsection{Event Types}

\api{EventType} values mirror Qt's event type numbers. Most applications handle these indirectly through higher-level callbacks, but the numeric values are useful when inspecting generic events, logs, or test helpers.

{\small
\begin{longtable}{@{}p{0.38\textwidth}p{0.10\textwidth}p{0.44\textwidth}@{}}
\toprule
Name & Value & Typical use \\
\midrule
\api{None} & 0 & No event or an unspecified event. \\
\api{Timer} & 1 & Timer delivery. \\
\api{MouseButtonPress} & 2 & Mouse or pointer button press. \\
\api{MouseButtonRelease} & 3 & Mouse or pointer button release. \\
\api{MouseButtonDblClick} & 4 & Mouse double-click. \\
\api{MouseMove} & 5 & Mouse or pointer movement. \\
\api{KeyPress} & 6 & Keyboard press. \\
\api{KeyRelease} & 7 & Keyboard release. \\
\api{FocusIn} & 8 & Widget receives keyboard focus. \\
\api{FocusOut} & 9 & Widget loses keyboard focus. \\
\api{Enter} & 10 & Pointer enters a widget. \\
\api{Leave} & 11 & Pointer leaves a widget. \\
\api{Paint} & 12 & Widget repaint request. \\
\api{Move} & 13 & Widget position changes. \\
\api{Resize} & 14 & Widget size changes. \\
\api{Create} & 15 & Native widget creation. \\
\api{Destroy} & 16 & Native widget destruction. \\
\api{Show} & 17 & Widget becomes visible. \\
\api{Hide} & 18 & Widget becomes hidden. \\
\api{Close} & 19 & Window close request. \\
\api{Quit} & 20 & Application quit request. \\
\api{ParentChange} & 21 & Object or widget parent changes. \\
\api{ThreadChange} & 22 & Object affinity moves to another thread. \\
\api{FocusAboutToChange} & 23 & Focus is about to move. \\
\api{WindowActivate} & 24 & Window becomes active. \\
\api{WindowDeactivate} & 25 & Window becomes inactive. \\
\api{ShowToParent} & 26 & Child widget is shown to its parent. \\
\api{HideToParent} & 27 & Child widget is hidden from its parent. \\
\api{Wheel} & 31 & Mouse wheel or trackpad scroll. \\
\api{WindowTitleChange} & 33 & Window title changes. \\
\api{WindowIconChange} & 34 & Window icon changes. \\
\api{ApplicationWindowIconChange} & 35 & Application window icon changes. \\
\api{ApplicationFontChange} & 36 & Application font changes. \\
\api{ApplicationLayoutDirectionChange} & 37 & Application layout direction changes. \\
\api{ApplicationPaletteChange} & 38 & Application palette changes. \\
\api{PaletteChange} & 39 & Widget palette changes. \\
\api{Clipboard} & 40 & Clipboard contents change. \\
\api{Speech} & 42 & Platform speech event. \\
\api{MetaCall} & 43 & Queued signal, slot, or meta-call delivery. \\
\api{SockAct} & 50 & Socket notifier activity. \\
\api{ShortcutOverride} & 51 & Shortcut override opportunity before key handling. \\
\api{DeferredDelete} & 52 & Deferred object deletion. \\
\api{DragEnter} & 60 & Drag operation enters a widget. \\
\api{DragMove} & 61 & Drag operation moves over a widget. \\
\api{DragLeave} & 62 & Drag operation leaves a widget. \\
\api{Drop} & 63 & Drop operation completes. \\
\bottomrule
\end{longtable}
}

\subsection{System Tray Enums}

The tray wrapper uses two small enums: one for how the user activated the tray icon, and one for the icon style shown in a tray message.

{\small
\begin{longtable}{@{}p{0.38\textwidth}p{0.10\textwidth}p{0.44\textwidth}@{}}
\toprule
Name & Value & Typical use \\
\midrule
\api{SystemTrayActivationReason::Unknown} & 0 & The platform did not report a more specific activation reason. \\
\api{SystemTrayActivationReason::Context} & 1 & The user requested the context menu, often by secondary click. \\
\api{SystemTrayActivationReason::DoubleClick} & 2 & The tray icon was double-clicked. \\
\api{SystemTrayActivationReason::Trigger} & 3 & The primary tray activation, often a single click. \\
\api{SystemTrayActivationReason::MiddleClick} & 4 & The tray icon was activated with the middle mouse button. \\
\api{SystemTrayMessageIcon::NoIcon} & 0 & Show a message without a status icon. \\
\api{SystemTrayMessageIcon::Information} & 1 & Use an informational status icon. \\
\api{SystemTrayMessageIcon::Warning} & 2 & Use a warning status icon. \\
\api{SystemTrayMessageIcon::Critical} & 3 & Use a critical or error status icon. \\
\bottomrule
\end{longtable}
}

\subsection{Common UI Enums}

{\small
\begin{longtable}{@{}p{0.34\textwidth}p{0.56\textwidth}@{}}
\toprule
Enum & Values \\
\midrule
\api{DialogCode} & \begin{tabular}[t]{@{}l@{}}\api{Rejected = 0}\\\api{Accepted = 1}\end{tabular} \\
\api{Orientation} & \begin{tabular}[t]{@{}l@{}}\api{Horizontal = 1}\\\api{Vertical = 2}\end{tabular} \\
\api{SortOrder} & \begin{tabular}[t]{@{}l@{}}\api{Ascending = 0}\\\api{Descending = 1}\end{tabular} \\
\api{CheckState} & \begin{tabular}[t]{@{}l@{}}\api{Unchecked = 0}\\\api{PartiallyChecked = 1}\\\api{Checked = 2}\end{tabular} \\
\api{FileDialogAcceptMode} & \begin{tabular}[t]{@{}l@{}}\api{Open = 0}\\\api{Save = 1}\end{tabular} \\
\api{FileDialogFileMode} & \begin{tabular}[t]{@{}l@{}}\api{AnyFile = 0}\\\api{ExistingFile = 1}\\\api{Directory = 2}\\\api{ExistingFiles = 3}\end{tabular} \\
\bottomrule
\end{longtable}
}

\subsection{Model/View Enums and Flags}

{\small
\begin{longtable}{@{}p{0.34\textwidth}p{0.56\textwidth}@{}}
\toprule
Enum & Values \\
\midrule
\api{ItemDataRole} & \begin{tabular}[t]{@{}l@{}}\api{Display = 0}\\\api{Decoration = 1}\\\api{Edit = 2}\\\api{ToolTip = 3}\\\api{StatusTip = 4}\\\api{WhatsThis = 5}\\\api{Font = 6}\\\api{TextAlignment = 7}\\\api{Background = 8}\\\api{Foreground = 9}\\\api{CheckState = 10}\\\api{AccessibleText = 11}\\\api{AccessibleDescription = 12}\\\api{SizeHint = 13}\\\api{InitialSortOrder = 14}\\\api{User = 256}\end{tabular} \\
\api{ItemFlag} & \begin{tabular}[t]{@{}l@{}}\api{None = 0}\\\api{Selectable = 1}\\\api{Editable = 2}\\\api{DragEnabled = 4}\\\api{DropEnabled = 8}\\\api{UserCheckable = 16}\\\api{Enabled = 32}\\\api{AutoTristate = 64}\\\api{NeverHasChildren = 128}\\\api{UserTristate = 256}\end{tabular} \\
\api{ItemSelectionMode} & \begin{tabular}[t]{@{}l@{}}\api{NoSelection = 0}\\\api{SingleSelection = 1}\\\api{MultiSelection = 2}\\\api{ExtendedSelection = 3}\\\api{ContiguousSelection = 4}\end{tabular} \\
\api{ItemSelectionBehavior} & \begin{tabular}[t]{@{}l@{}}\api{SelectItems = 0}\\\api{SelectRows = 1}\\\api{SelectColumns = 2}\end{tabular} \\
\api{SelectionFlag} & \begin{tabular}[t]{@{}l@{}}\api{NoUpdate = 0}\\\api{Clear = 1}\\\api{Select = 2}\\\api{Deselect = 4}\\\api{Toggle = 8}\\\api{Current = 16}\\\api{Rows = 32}\\\api{Columns = 64}\\\api{SelectCurrent = 18}\\\api{ToggleCurrent = 24}\\\api{ClearAndSelect = 3}\end{tabular} \\
\api{DropAction} & \begin{tabular}[t]{@{}l@{}}\api{IgnoreAction = 0}\\\api{CopyAction = 1}\\\api{MoveAction = 2}\\\api{LinkAction = 4}\end{tabular} \\
\api{EditTrigger} & \begin{tabular}[t]{@{}l@{}}\api{NoEditTriggers = 0}\\\api{CurrentChanged = 1}\\\api{DoubleClicked = 2}\\\api{SelectedClicked = 4}\\\api{EditKeyPressed = 8}\\\api{AnyKeyPressed = 16}\\\api{AllEditTriggers = 31}\end{tabular} \\
\bottomrule
\end{longtable}
}

\subsection{Graphics View Enums and Flags}

{\small
\begin{longtable}{@{}p{0.34\textwidth}p{0.56\textwidth}@{}}
\toprule
Enum & Values \\
\midrule
\api{GraphicsItemFlag} & \begin{tabular}[t]{@{}l@{}}\api{None = 0}\\\api{ItemIsMovable = 1}\\\api{ItemIsSelectable = 2}\\\api{ItemIsFocusable = 4}\\\api{ItemClipsToShape = 8}\\\api{ItemClipsChildrenToShape = 16}\\\api{ItemIgnoresTransformations = 32}\\\api{ItemIgnoresParentOpacity = 64}\\\api{ItemDoesntPropagateOpacityToChildren = 128}\\\api{ItemStacksBehindParent = 256}\\\api{ItemUsesExtendedStyleOption = 512}\\\api{ItemHasNoContents = 1024}\\\api{ItemSendsGeometryChanges = 2048}\\\api{ItemAcceptsInputMethod = 4096}\\\api{ItemNegativeZStacksBehindParent = 8192}\\\api{ItemIsPanel = 16384}\\\api{ItemSendsScenePositionChanges = 65536}\end{tabular} \\
\api{GraphicsItemCacheMode} & \begin{tabular}[t]{@{}l@{}}\api{NoCache = 0}\\\api{ItemCoordinateCache = 1}\\\api{DeviceCoordinateCache = 2}\end{tabular} \\
\api{GraphicsSceneItemIndexMethod} & \begin{tabular}[t]{@{}l@{}}\api{BspTreeIndex = 0}\\\api{NoIndex = -1}\end{tabular} \\
\api{GraphicsSceneLayer} & \begin{tabular}[t]{@{}l@{}}\api{ItemLayer = 1}\\\api{BackgroundLayer = 2}\\\api{ForegroundLayer = 4}\\\api{AllLayers = 65535}\end{tabular} \\
\api{GraphicsItemSelectionMode} & \begin{tabular}[t]{@{}l@{}}\api{IntersectsItemShape = 0}\\\api{ContainsItemShape = 1}\\\api{IntersectsItemBoundingRect = 2}\\\api{ContainsItemBoundingRect = 3}\end{tabular} \\
\api{GraphicsViewDragMode} & \begin{tabular}[t]{@{}l@{}}\api{NoDrag = 0}\\\api{ScrollHandDrag = 1}\\\api{RubberBandDrag = 2}\end{tabular} \\
\api{GraphicsViewViewportAnchor} & \begin{tabular}[t]{@{}l@{}}\api{NoAnchor = 0}\\\api{AnchorViewCenter = 1}\\\api{AnchorUnderMouse = 2}\end{tabular} \\
\api{GraphicsViewViewportUpdateMode} & \begin{tabular}[t]{@{}l@{}}\api{FullViewportUpdate = 0}\\\api{MinimalViewportUpdate = 1}\\\api{SmartViewportUpdate = 2}\\\api{NoViewportUpdate = 3}\\\api{BoundingRectViewportUpdate = 4}\end{tabular} \\
\api{GraphicsViewOptimizationFlag} & \begin{tabular}[t]{@{}l@{}}\api{DontSavePainterState = 1}\\\api{DontAdjustForAntialiasing = 2}\\\api{IndirectPainting = 4}\end{tabular} \\
\api{GraphicsViewCacheModeFlag} & \begin{tabular}[t]{@{}l@{}}\api{CacheNone = 0}\\\api{CacheBackground = 1}\end{tabular} \\
\api{GraphicsPixmapItemShapeMode} & \begin{tabular}[t]{@{}l@{}}\api{MaskShape = 0}\\\api{BoundingRectShape = 1}\\\api{HeuristicMaskShape = 2}\end{tabular} \\
\api{GraphicsSceneContextMenuEventReason} & \begin{tabular}[t]{@{}l@{}}\api{Mouse = 0}\\\api{Keyboard = 1}\\\api{Other = 2}\end{tabular} \\
\bottomrule
\end{longtable}
}

\subsection{Rendering and Image Enums}

{\small
\begin{longtable}{@{}p{0.34\textwidth}p{0.56\textwidth}@{}}
\toprule
Enum & Values \\
\midrule
\api{AspectRatioMode} & \begin{tabular}[t]{@{}l@{}}\api{Ignore = 0}\\\api{Keep = 1}\\\api{KeepByExpanding = 2}\end{tabular} \\
\api{TransformationMode} & \begin{tabular}[t]{@{}l@{}}\api{Fast = 0}\\\api{Smooth = 1}\end{tabular} \\
\api{PenStyle} & \begin{tabular}[t]{@{}l@{}}\api{NoPen = 0}\\\api{SolidLine = 1}\\\api{DashLine = 2}\\\api{DotLine = 3}\\\api{DashDotLine = 4}\\\api{DashDotDotLine = 5}\\\api{CustomDashLine = 6}\end{tabular} \\
\api{PenCapStyle} & \begin{tabular}[t]{@{}l@{}}\api{FlatCap = 0}\\\api{SquareCap = 16}\\\api{RoundCap = 32}\end{tabular} \\
\api{PenJoinStyle} & \begin{tabular}[t]{@{}l@{}}\api{MiterJoin = 0}\\\api{BevelJoin = 64}\\\api{RoundJoin = 128}\\\api{SvgMiterJoin = 256}\end{tabular} \\
\api{PainterCompositionMode} & \begin{tabular}[t]{@{}l@{}}\api{SourceOver = 0}\\\api{DestinationOver = 1}\\\api{Clear = 2}\\\api{Source = 3}\\\api{Destination = 4}\\\api{SourceIn = 5}\\\api{DestinationIn = 6}\\\api{SourceOut = 7}\\\api{DestinationOut = 8}\\\api{SourceAtop = 9}\\\api{DestinationAtop = 10}\\\api{Xor = 11}\\\api{Plus = 12}\\\api{Multiply = 13}\\\api{Screen = 14}\\\api{Overlay = 15}\\\api{Darken = 16}\\\api{Lighten = 17}\\\api{ColorDodge = 18}\\\api{ColorBurn = 19}\\\api{HardLight = 20}\\\api{SoftLight = 21}\\\api{Difference = 22}\\\api{Exclusion = 23}\end{tabular} \\
\bottomrule
\end{longtable}
}

\section{Keeping This Appendix Current}

When a new wrapper enum becomes visible in examples, callbacks, or public setters, add it here if a reader would otherwise need to jump into the source to understand a small code sample. Internal or rarely used enums can stay in the API source until they become part of the guide's narrative.
